\section{Conclusion}

Cette étude avait pour objectif d'identifier des structures thématiques récurrentes dans l'œuvre romanesque de Zola à partir de méthodes non supervisées. Le corpus a d'abord été nettoyé, segmenté en paquets de phrases, puis lemmatisé afin de produire des unités textuelles comparables et exploitables par les modèles.

L'application de le NMF a permis de faire apparaître des topics interprétables, associés à des dimensions importantes de l'univers de zola : le travail, la mine, la religion, la famille, l'argent, la guerre, le commerce, la ville ou encore l'espace domestique. Le recours à K-means a ensuite permis de regrouper les segments selon leur proximité lexicale et d'observer des familles de motifs complémentaires.

La comparaison des deux méthodes montre que certains ensembles thématiques apparaissent de manière stable, malgré des fonctionnements algorithmiques différents. le NMF offre une lecture plus graduelle des thèmes, tandis que K-means produit une classification plus nette des segments. Leur combinaison permet donc de renforcer l'analyse en croisant deux formes de structuration du corpus.

Ces résultats doivent toutefois être interprétés avec prudence. Ils dépendent des choix de nettoyage, de lemmatisation, de segmentation, de vectorisation et du nombre de topics ou de clusters retenus. L'analyse proposée ne remplace donc pas une lecture littéraire des textes, mais elle fournit un outil exploratoire permettant de repérer des régularités lexicales et thématiques à l'échelle d'un corpus volumineux.
